% !TeX root = ../main.tex
\chapter{Euclidean Space}
\section{Norm and Inner Product}
If we call something Euclidean, we need to define so called length and angle, so we introduce the concepts of norm and inner product on a linear space.
\begin{definition}{Norm}
  we define $\phi:V \to \mathbb{R}$ s.t. $\forall x \in V$, denote $\left\| x \right\|$:
  \begin{enumerate}
    \item $\left\| x \right\| \geq 0$, $\left\| x \right\| = 0 \leftrightarrow x = O$ 
    \item $\left\| c x \right\| = |c|\left\| x \right\|,\forall x \in V, c \in \mathbb{R}$
    \item $\left\| x + y\right\| \leq \left\| x \right\| + \left\| y \right\|$
  \end{enumerate}
\end{definition}
\begin{definition}{Inner Product}
  define $\psi :V \circ V\to \mathbb{R}, \forall x, y,z \in V$, we denote $(x,y)$:
  \begin{enumerate}
    \item $(x,y)=(y,x)$
    \item $(x,x) \geq 0, (x,x) = 0 \leftrightarrow x = O$
    \item $(cx,y) = c(x,y)$
    \item $(x +y,z) = (x,z)+(y,z)$
  \end{enumerate}
\end{definition}
A Real Linear Space with an Inner product is called an Euclidean Space, and every inner product comes with an induced norm,
$\left\| x \right\| = \sqrt{(x,x)}$
and we can obtain the Cauchy-Schwartz Inequality
\[|(x,y)| \leq \left\| x \right\|\left\| y \right\|\]
\section{Orthogonality}


